\documentclass[11pt,a4paper]{article}
\usepackage[T2A]{fontenc}
\usepackage[utf8]{inputenc}
\usepackage[russian,english]{babel}
\usepackage{amsmath,amssymb,amsthm,mathtools}
\usepackage{setspace}
\usepackage{enumitem}
\usepackage[hidelinks]{hyperref}
\onehalfspacing
\renewcommand{\leq}{\leqslant}
\renewcommand{\le}{\leqslant}
\renewcommand{\geq}{\geqslant}
\renewcommand{\ge}{\geqslant}
\newcommand{\head}[1]{\par\bigskip\noindent\textbf{\underline{#1}}\quad}
\begin{document}
\begin{center}
{\huge\bfseries Methods of proof -- 1\par}\vspace{1.2em}
{\Large The pigeonhole principle: discrete, continuous, probabilistic\par}\vspace{0.8em}
{\foreignlanguage{russian}{13 сентября 2026 года}\par}
\end{center}\vspace{0.5em}

\noindent\textbf{\underline{Theoretical minimum.}}\par

\head{Fact 1 (the discrete form).} If $n$ objects are put into $k$ boxes, some box contains at
least $\lceil n/k\rceil$ of them. The whole art is choosing the boxes: residues, dyadic scales,
sign patterns, first digits, positions in an ordering.

\head{Theorem 1 (the continuous form).} If sets $A_1,\dots,A_n$ inside a set of measure $\mu$
have total measure greater than $k\mu$, some point lies in at least $k+1$ of them (Blichfeldt).
Two consequences: Minkowski's theorem --- a centrally symmetric convex body of volume greater
than $2^n$ contains a non-zero lattice point (Convex -- 10) --- and Steinhaus' theorem: if $A$
has positive measure, then $A-A$ contains an interval around zero.

\head{Theorem 2 (Dirichlet's approximation).} For every irrational $\alpha$ there are
infinitely many fractions with $\bigl|\alpha-\frac pq\bigr|<\frac1{q^2}$; the proof is the
pigeonhole applied to the fractional parts $\{q\alpha\}$ in $N$ boxes. The same argument gives
the density of $\{n\alpha\}$, which was upgraded to equidistribution in Trig-5.

\head{Fact 2 (the probabilistic form).} There is always an element at least as good as the
average: if $\mathrm E\,X=m$ then $X\ge m$ with positive probability. This is the pigeonhole
with weights, and the second moment refines it --- a small variance means that almost every
element is close to the average (Probability -- 1 and Inequalities -- 8).

\bigskip\noindent\rule{\linewidth}{0.4pt}

\begin{enumerate}[leftmargin=2.2em,itemsep=0.6em]\interlinepenalty=10000
\item Prove the discrete pigeonhole in the form with $\lceil n/k\rceil$ and deduce Dirichlet's
approximation theorem; then prove that for every irrational $\alpha$ and every $N$ there is
$q\le N$ with $\|q\alpha\|<\frac1N$, and deduce the density of $\{n\alpha\}$ in $[0,1]$.

\item Prove Blichfeldt's lemma and derive from it both Minkowski's theorem on a symmetric convex
body and Steinhaus' theorem that $A-A$ contains an interval whenever $m(A)>0$. Deduce that an
additive function bounded on a set of positive measure is linear (compare with Continuous
functions -- 6).

\item (IMC) Let $X$ be a set of $\binom{2k-4}{k-2}+1$ real numbers, $k\ge2$. Prove that there
exists a monotone sequence $\{x_i\}_{i=1}^k\subseteq X$ such that $|x_{i+1}-x_1|\ge2|x_i-x_1|$
for all $i=2,\dots,k-1$.

\item Prove that among any $n+1$ numbers chosen from $\{1,2,\dots,2n\}$ some one divides
another, and that among any $n^2+1$ distinct reals there is a monotone subsequence of length
$n+1$; identify the boxes in each case.

\item Prove by a counting argument that for $k\ge3$ and $n<2^{k/2}$ there exists a
two-colouring of the edges of $K_n$ with no monochromatic $K_k$, and state the general
principle: to prove existence, count the bad configurations and show that they cannot cover
everything.
\end{enumerate}
\end{document}
